\documentclass[aspectratio=169,12pt]{beamer}

% ==================================================
% PAQUETES
% ==================================================
\usepackage[spanish,es-noshorthands]{babel}
\usepackage[utf8]{inputenc}
\usepackage[T1]{fontenc}
\usepackage{lmodern}
\usepackage{amsmath,amssymb}
\usepackage{tikz}
\usepackage{xcolor}

\usetheme{Madrid}
\renewcommand{\familydefault}{\sfdefault}
\setbeamertemplate{navigation symbols}{}

% ==================================================
% COLORES
% ==================================================
\definecolor{Primary}{HTML}{0B3C5D}
\definecolor{Accent}{HTML}{2E8B57}
\definecolor{Soft}{HTML}{EEF4F8}
\definecolor{Dark}{HTML}{1F2937}

\setbeamercolor{background canvas}{bg=white}
\setbeamercolor{normal text}{fg=Dark}
\setbeamercolor{frametitle}{bg=Soft, fg=Primary}
\setbeamercolor{block title}{bg=Primary, fg=white}
\setbeamercolor{block body}{bg=Primary!5}

% ==================================================
% COMANDOS
% ==================================================
\newcommand{\IdeaClave}[1]{
	\begin{block}{Idea clave}
		#1
	\end{block}
}

\newcommand{\Paso}[1]{
	\textbf{Paso:} #1
}

% ==================================================
\title{Factorización}
\subtitle{Clase develada paso a paso}
\author{Luis Tulio González}
\date{2026}

% ==================================================
\begin{document}
	
	% ==================================================
	% PORTADA PRO
	% ==================================================
	\begin{frame}[plain]
		
		\begin{tikzpicture}[remember picture,overlay]
			
			% FONDO SUPERIOR
			\fill[Primary] (current page.north west) rectangle ([yshift=-3cm]current page.north east);
			
			% LINEA DE ACENTO
			\fill[Accent] ([yshift=-3cm]current page.north west) rectangle ([yshift=-3.2cm]current page.north east);
			
			% TITULO
			\node[anchor=north west, text=white] 
			at ([xshift=1cm,yshift=-1cm]current page.north west){
				\begin{minipage}{0.85\paperwidth}
					{\fontsize{23}{25}\selectfont\bfseries Fundamentos de Matemática Básica}\\[0.4cm]
					{\large Unidad I — Clase: Factorización}\\[0.2cm]					
				\end{minipage}
			};
			
		\end{tikzpicture}
		
		\vspace{4.5cm}
		
		% DATOS
		\begin{center}
			{\large\bfseries Luis Tulio González Alcaino}\\[0.3cm]
			{\normalsize Universidad Santo Tomás}\\
			{\normalsize Departamento Ciencias Básicas — Talca}\\[0.5cm]
			{\normalsize Semestre I — 2026}
		\end{center}
		
	\end{frame}
	
	% ==================================================
	\begin{frame}{Motivación}
		
		\IdeaClave{
			En Agronomía, muchas expresiones algebraicas aparecen al modelar:
			\begin{itemize}
				\item costos de producción,
				\item dosis de fertilización,
				\item áreas de cultivo,
				\item rendimiento por parcela,
				\item variaciones de riego.
			\end{itemize}
		}
		
		\pause
		
		\begin{block}{Objetivo}
			Reconocer el tipo de factorización adecuado, interpretar patrones y evitar errores.
		\end{block}
		
	\end{frame}
	
	% ==================================================
	\begin{frame}{¿Qué significa factorizar?}
		
		\begin{block}{Definición}
			Factorizar es el proceso inverso de multiplicar.
		\end{block}
		
		\pause
		
		\begin{exampleblock}{Idea}
			Si desarrollas obtienes suma,\\
			si factorizas vuelves a producto.
		\end{exampleblock}
		
		\pause
		
		\begin{block}{Ejemplo guiado}
			\[
			6x + 12
			\]
			\pause
			
			\Paso{Buscar factor común}
			
			\pause
			\[
			6x + 12 = 6(x+2)
			\]
			
		\end{block}
		
	\end{frame}
	
	% ==================================================
	\begin{frame}{Interpretación agronómica}
		
		\begin{exampleblock}{Situación}
			Costo total:
			\[
			6x + 12
			\]
			\pause
			
			\Paso{Factorizamos}
			\[
			6(x+2)
			\]
			
			\pause
			
			\Paso{Interpretamos}
			\begin{itemize}
				\item 6 = costo por unidad
				\item x+2 = cantidad total
			\end{itemize}
			
		\end{exampleblock}
		
	\end{frame}
	
	% ==================================================
	\section{Factor común}
	
	\begin{frame}{Factor común}
		
		\begin{block}{Patrón}
			\[
			ab + ac = a(b+c)
			\]
		\end{block}
		
		\pause
		
		\begin{block}{Ejemplo}
			\[
			8x + 16
			\]
			
			\pause
			\Paso{Factor común = 8}
			
			\pause
			\[
			8(x+2)
			\]
			
		\end{block}
		
	\end{frame}
	
	% ==================================================
	\section{Diferencia de cuadrados}
	
	\begin{frame}{Suma por su diferencia}
		
		\begin{block}{Patrón}
			\[
			a^2 - b^2 = (a-b)(a+b)
			\]
		\end{block}
		
		\pause
		
		\begin{block}{Ejemplo}
			\[
			x^2 - 25
			\]
			
			\pause
			\Paso{Identificar cuadrados}
			\[
			x^2 - 5^2
			\]
			
			\pause
			\[
			(x-5)(x+5)
			\]
			
		\end{block}
		
	\end{frame}
	
	% ==================================================
	\section{Trinomio cuadrado perfecto}
	
	\begin{frame}{Trinomio cuadrado perfecto}
		
		\begin{block}{Patrón}
			\[
			a^2 + 2ab + b^2 = (a+b)^2
			\]
		\end{block}
		
		\pause
		
		\begin{block}{Ejemplo}
			\[
			x^2 + 6x + 9
			\]
			
			\pause
			\Paso{Verificar patrón}
			
			\pause
			\[
			(x+3)^2
			\]
			
		\end{block}
		
	\end{frame}
	
	% ==================================================
	\section{Trinomio x²+bx+c}
	
	\begin{frame}{Trinomio \(x^2+bx+c\)}
		
		\begin{block}{Idea clave}
			Buscar dos números $m$ y $n$ tal que:
			\begin{itemize}
				\item multipliquen \(m \cdot n= c\)
				\item sumen \(m+n=b\)
				\item \(x^2+bx+c=(x+m)(x+n)\)
			\end{itemize}
		\end{block}
		
		\pause
		
		\begin{block}{Ejemplo}
			\[
			x^2 + 7x + 12
			\]
			
			\pause
			\Paso{Buscar números}
			
			\pause
			\[
			\quad 3\cdot4 = 12
			\quad ; 3 + 4 = 7
			\]
			
			\pause
			\[
			(x+3)(x+4)
			\]
			
		\end{block}
		
	\end{frame}
	
	% ==================================================
	\section{Errores}
	
	\begin{frame}{Errores frecuentes}
		
		\begin{alertblock}{Error}
			\[
			a^2 + b^2 \neq (a+b)^2
			\]
		\end{alertblock}
		
		\pause
		
		\begin{alertblock}{Error}
			No todo trinomio es cuadrado perfecto
		\end{alertblock}
		
		\pause
		
		\begin{alertblock}{Error}
			Elegir mal los números en \(x^2+bx+c\)
		\end{alertblock}
		
	\end{frame}
	
	% ==================================================
	\section{Ejercicios}
	
	\begin{frame}{Ejercicios guiados}
		
		Factoriza:
		
		\begin{enumerate}
			\item \(6x + 18\)
			\pause
			\item \(x^2 - 36\)
			\pause
			\item \(x^2 + 10x + 25\)
			\pause
			\item \(x^2 + 9x + 20\)
		\end{enumerate}
		
	\end{frame}
	
	% ==================================================
	\begin{frame}{Cierre}
		
		\IdeaClave{
			Factorizar no es memorizar, es reconocer patrones.
		}
		
		\pause
		
		\begin{block}{Clave}
			\begin{itemize}
				\item Observa la forma
				\item Identifica el patrón
				\item Aplica la técnica correcta
			\end{itemize}
		\end{block}
		
	\end{frame}
	
	% ==================================================
\end{document}